\documentclass[a4paper,12pt,UTF8,fontset=none]{ctexart}
\usepackage{geometry} % 页面设置
\usepackage{xeCJK}
\usepackage{enumitem}
\usepackage{xcolor} % Color support for listings
\usepackage{graphicx}
\usepackage{float}
\usepackage{pdfpages}
\usepackage{amsmath, amsthm, amssymb}
\usepackage{listings}
\usepackage{longtable}
\usepackage{booktabs} % 添加booktabs宏包以支持三线表
\usepackage{tabularx} % 新增tabularx宏包用于自动调整列宽
\usepackage{array}      % 列格式控制
\usepackage{caption}    % 表格标题格式
\usepackage{placeins} 
\usepackage{newtxtext,newtxmath}
\usepackage{fontspec}
\usepackage{titlesec}
\usepackage{bm} % 加粗符号
\usepackage{background} % 添加水印宏包
\usepackage{hyperref} % 超链接引用
\usepackage{cleveref}
\newtheorem{theorem}{定理}[section] % 定理编号按章节
\newtheorem{proofstep}{步骤}[theorem] % 证明步骤编号依附定理

% 自定义关键项高亮命令
\newcommand{\highlight}[1]{\textcolor{red}{\boxed{#1}}}
\backgroundsetup{ % 水印参数设置
    scale=3,          % 水印文字大小（可调）
    angle=0,         % 水印旋转角度
    opacity=0.25,     % 透明度（0-1）
    contents={\includegraphics[width=0.5\textwidth]{static/pdfs/水印.pdf}} % 改为正确的水印文件
    % contents={lanshi}
}
% 设置英文主字体为 Times New Roman
\setmainfont{Times New Roman}[Path=D:/Program Files/MiKTeX/fonts/custom/, Extension=.otf]

% 设置英文粗体字体为 Times New Roman Bold
\newfontfamily\enboldfont{Times New Roman Bold}[Path=D:/Program Files/MiKTeX/fonts/custom/, Extension=.otf]

% 设置中文正文字体为宋体
\setCJKmainfont{SimSun}[Path=D:/Program Files/MiKTeX/fonts/custom/, Extension=.ttf]

% 设置中文黑体字体
\setCJKfamilyfont{zhhei}{SimHei}[Path=D:/Program Files/MiKTeX/fonts/custom/, Extension=.ttf]
\newcommand{\heiti}{\CJKfamily{zhhei}}

% 自定义编号格式
\renewcommand{\thesection}{\arabic{section}}
\renewcommand{\thesubsection}{\thesection.\arabic{subsection}}
\renewcommand{\thesubsubsection}{\thesubsection.\arabic{subsubsection}}
% 设置标题格式
\titleformat{\section}
  {\centering\heiti\zihao{-3}\enboldfont}{\thesection}{1em}{}
\titleformat{\subsection}
  {\heiti\zihao{4}\enboldfont}{\thesubsection}{1em}{}
\titleformat{\subsubsection}
  {\heiti\zihao{-4}\enboldfont}{\thesubsubsection}{1em}{}


% 定义计数器
\newcounter{mycounter}

% 定义命令用于显示编号
\newcommand{\myitem}{%
  \par \stepcounter{mycounter}%
  \textbf{\themycounter.}~%
}
% 在每个 section 开始时重置计数器
\let\oldsection\section
\renewcommand{\section}{%
  \setcounter{mycounter}{0}%
  \oldsection%
}
% 设置代码块样式
\lstset{
    basicstyle=\ttfamily,
    columns=fullflexible,
    frame=single,
    breaklines=true,
    postbreak=\mbox{\textcolor{red}{$\hookrightarrow$}\space},
    language=SQL,
    keywordstyle=\color{blue},
    commentstyle=\color{gray},    
    rulecolor=\color{black!30},%边框颜色
    stringstyle=\color{red},
    escapeinside={\%*}{*)},
    showstringspaces=false,
    captionpos=b % 设置标题位置, b表示在底部
}
\geometry{left=2.5cm,right=2.5cm,top=2.5cm,bottom=2.5cm} % 页边距

\begin{document}

% 制作封面页
\begin{titlepage}
    \centering
    \vspace*{\fill}
    {\LARGE\bfseries 人工智能学习笔记\par}
    \vspace{2cm}
    {\Huge\bfseries 西瓜书学习笔记\par}
    \vspace{2cm}
    {\Large 烂石\par}
    \vspace{1cm}
    {\large \today \par}
    \vspace{4cm}
    \includegraphics[width=0.5\textwidth]{static/images/logo.png}
    \vspace*{\fill}
    \thispagestyle{empty} % 封面不显示页码
    \newpage
\end{titlepage}
\section{绪论}
\subsection{引言}
略
\subsection{基本术语}
\subsubsection{数据集-data set}
训练集-training data
\begin{center}
    \begin{tikzpicture}[
        mindmap,
        concept color=orange!30,  % 中心节点颜色
        every node/.style=concept,
        level 1/.append style={
            sibling angle=120,
            level distance=5cm,
            concept color=blue!20
        },
        level 2/.append style={
            sibling angle=60,
            level distance=3.5cm,
            concept color=green!20
        },
        level 3/.append style={
            sibling angle=45,
            level distance=2.5cm,
            concept color=red!20
        }
    ]
    
    % 中心节点
    \node {机器学习类型} 
        % 第一层级分支
        child [grow=30] { node {监督学习}
            child { node {分类}
                child { node {二分类}
                    child { node {逻辑回归} }
                    child { node {SVM} }
                }
                child { node {多分类}
                    child { node {Softmax回归} }
                    child { node {决策树} }
                }
            }
            child { node {回归}
                child { node {线性回归} }
                child { node {多项式回归} }
                child { node {Lasso回归} }
            }
        }
        child [grow=150] { node {无监督学习}
            child { node {聚类}
                child { node {K-Means} }
                child { node {层次聚类} }
                child { node {DBSCAN} }
            }
            child { node {降维}
                child { node {PCA} }
                child { node {t-SNE} }
            }
        }
        child [grow=270] { node {半监督学习}
            child { node {标签传播} }
            child { node {自训练} }
        };
    
    \end{tikzpicture}
    \end{center}

\end{document}